\section{GAN: Redes generativas adversarias}

\subsection{Del ataque adversario a los modelos generativos}

La idea detrás de las GANs (Generative Adversarial Networks) se origina en el ataque adversario. El docente ilustró esto con un ejemplo clásico: un clasificador entrenado puede identificar correctamente un panda, pero si se añade una perturbación microscópica casi imperceptible para el ojo humano (ruido) a la imagen, el clasificador la clasificará erróneamente como un avestruz —incluso si se amplía la diferencia 10 veces, es difícil que el ojo humano distinga la diferencia.

\begin{knowledgebox}{Del ataque adversario al salto conceptual de las GANs}
El ataque adversario reveló una verdad: las redes neuronales pueden ser extremadamente sensibles a cambios sutiles imperceptibles para el ojo humano. Ian Goodfellow, el creador de las GANs, se inspiró en esto: si una red (atacante/generador) puede generar muestras que engañen a otra red (defensor/discriminador), entonces mediante la competencia entre ambas, el generador eventualmente aprenderá a producir datos realistas.
\end{knowledgebox}

\subsection{Arquitectura de las GANs}

Las GANs están compuestas por dos redes neuronales:

\begin{itemize}
    \item \textbf{Generator (Generador)} $G_\theta$: recibe muestras latentes $z \sim p(z)$ y produce puntos de datos ``falsos'' $G_\theta(z)$.
    \item \textbf{Discriminator (Discriminador)} $D_\phi$: recibe un punto de datos y devuelve la probabilidad de que ese punto sea ``datos reales'', $D_\phi(x) \in {[}0, 1{]}$.
\end{itemize}

El generador y el discriminador compiten entre sí:
\begin{itemize}
    \item Objetivo del discriminador: distinguir correctamente los datos reales de los generados.
    \item Objetivo del generador: producir muestras que engañen al discriminador.
\end{itemize}

\subsection{Función de pérdida de las GANs: Juego Minimax}

El objetivo de entrenamiento de las GANs es un problema de optimización minimax:

$$\min_\theta \max_\phi \; V(G_\theta, D_\phi) = \mathbb{E}_{x \sim p_{\text{data}}} {[}\log D_\phi(x){]} + \mathbb{E}_{z \sim p(z)} {[}\log(1 - D_\phi(G_\theta(z))){]}$$

\begin{importantbox}{Significado de cada término en la función de pérdida de las GANs}
\begin{itemize}
    \item $\mathbb{E}_{x \sim p_{\text{data}}} {[}\log D_\phi(x){]}$: para los datos reales, el discriminador debe devolver una probabilidad alta (cercana a 1), haciendo que $\log D_\phi(x)$ sea lo más grande posible.
    \item $\mathbb{E}_{z \sim p(z)} {[}\log(1 - D_\phi(G_\theta(z))){]}$: para los datos generados $G_\theta(z)$, el discriminador debe devolver una probabilidad baja (cercana a 0), haciendo que $1 - D_\phi(G_\theta(z))$ se acerque a 1.
\end{itemize}
El \textbf{discriminador} mejora su capacidad de distinción \textbf{maximizando} $V$; el \textbf{generador} mejora su capacidad de engaño \textbf{minimizando} $V$.
\end{importantbox}

\subsubsection{Solución óptima del discriminador}

Manteniendo fijo el generador $G_\theta$, buscamos la solución óptima para el discriminador:

$$D^*_\phi(x) = \frac{p_{\text{data}}(x)}{p_{\text{data}}(x) + p_G(x)}$$

Donde $p_G(x)$ es la distribución inducida por el generador. Cuando $p_G = p_{\text{data}}$, entonces $D^*(x) = 1/2$, lo que significa que el discriminador no puede distinguir entre verdadero y falso —esto es precisamente el estado ideal hacia el cual converge el entrenamiento de las GANs.

\subsubsection{GANs con solución óptima global}

Se puede demostrar que cuando $p_G = p_{\text{data}}$, la función de pérdida alcanza su valor mínimo global $-2\log 2$. En este punto, la distribución del generador coincide completamente con la distribución de los datos.

\subsection{Dificultades en el entrenamiento de las GANs}

\begin{warningbox}{Problemas comunes en el entrenamiento de las GANs}
El problema de optimización minimax es \textbf{extremadamente difícil de optimizar} en la práctica y tiende a caer en óptimos locales pobres. Los modos de falla más frecuentes incluyen:
\begin{itemize}
    \item \textbf{Discriminador dominante (Discriminator Dominance)}: el discriminador converge mucho más rápido que el generador, lo que resulta en una señal de gradiente extremadamente débil (casi nula) para el generador, impidiendo un aprendizaje efectivo. Este es el modo de falla más común.
    \item \textbf{Colapso de modos (Mode Collapse)}: incluso si el entrenamiento progresa bien, el generador podría aprender a producir solo una parte de los modos de la distribución de datos, perdiendo así la diversidad de los datos. Por ejemplo, en datos multimodales, el generador podría concentrarse produciendo muestras de un solo modo.
    \item \textbf{Inestabilidad en el entrenamiento}: la curva de pérdida oscila violentamente, haciendo difícil determinar si se ha alcanzado la convergencia.
\end{itemize}
\end{warningbox}

El docente mencionó una analogía vívida: si estableces criterios de éxito extremadamente estrictos desde el principio (un discriminador demasiado fuerte), no aprenderás nada —esto es exactamente como las lecciones de la vida.

\subsubsection{Práctica ingenieril en el entrenamiento de las GANs}

En la práctica, entrenar GANs requiere numerosas técnicas de ingeniería:

\begin{itemize}
    \item Por lo general, se preentrena al generador durante unos pasos para que pueda producir salidas con algún significado antes de introducir el discriminador para el entrenamiento adversarial.
    \item Se necesita equilibrar cuidadosamente las tasas de aprendizaje y la proporción de pasos de entrenamiento de ambas redes.
    \item Muchas empresas dedicaron grandes equipos de ingeniería especializados a ajustar los parámetros de entrenamiento de las GANs antes del año 2020.
\end{itemize}

\subsection{Evolución de las GANs}

Desde que Goodfellow y sus colegas propusieron las GANs en 2013/2014, han surgido numerosas variantes. El docente mencionó brevemente varios hitos fundamentales:

\begin{table}[H]
\centering
\begin{tabular}{lll}
\toprule
Modelo & Año & Contribución principal \\
\midrule
GAN & 2014 & Propuesta del marco de entrenamiento adversario \\
DCGAN & 2015 & Introducción de la arquitectura convolucional \\
WGAN & 2017 & Sustitución de la divergencia JS por la distancia Wasserstein \\
Progressive GAN & 2017 & Estrategia de crecimiento progresivo \\
StyleGAN / StyleGAN2 & 2019/2020 & Generación de rostros de alta calidad \\
\bottomrule
\end{tabular}
\caption{Variantes importantes de las GANs (parcial)}
\end{table}

\begin{knowledgebox}{¿Por qué se pasó de las GANs a los modelos de difusión?}
El docente señaló que la razón principal por la cual la industria pasó de las GANs a los modelos de Difusión/Flujo no fue la calidad —StyleGAN ya podía generar imágenes de rostros extremadamente realistas en 2019—. El verdadero impulso fue: (1) \textbf{Dificultad en el entrenamiento} —la optimización minimax es inestable y requiere una gran inversión de ingeniería; (2) \textbf{Falta de diversidad} —el problema del colapso de modos hace que las muestras generadas carezcan de variedad. El modelo de difusión evita astutamente estos dos problemas al descomponer el proceso de generación en múltiples pasos de eliminación de ruido.
\end{knowledgebox}

\subsection{Resumen de la sección}

Las GANs aprenden la distribución de datos mediante un juego adversario entre un generador y un discriminador. Su marco teórico es elegante (el juego minimax puede coincidir exactamente con la distribución de datos en condiciones ideales), pero en la práctica enfrenta problemas graves como la inestabilidad del entrenamiento y el colapso de modos. Estos desafíos impulsaron el desarrollo de alternativas como las VAEs y, más tarde, los modelos de difusión.

\newpage
